\mychapter{Referencial Teórico}{cap:tec}\lhead{REFERENCIAL TEÓRICO}

Este capítulo consolida os conceitos, modelos e arcabouços teóricos que fundamentam a dissertação. A exposição organiza-se em seis eixos integrados: (i)~caracterização quantitativa da demanda intermitente, com ênfase na taxonomia de Syntetos-Boylan; (ii)~políticas clássicas e meta-heurísticas de controle de inventário; (iii)~aprendizado por reforço aplicado a decisões sequenciais de reposição; (iv)~métricas de qualidade de previsão adequadas ao regime intermitente; (v)~validação temporal sem vazamento de informação; e (vi)~testes estatísticos não paramétricos para comparação de múltiplas políticas. A unificação desses eixos no pipeline Kedro implementado constitui a contribuição metodológica central da dissertação.

\section{Demanda Intermitente: Definição e Caracterização}

\subsection{Conceito e Relevância}

Demanda intermitente é caracterizada por ciclos frequentes de demanda nula intercalados com valores positivos, muitas vezes de magnitude irregular. Esse padrão é prevalente em varejo regional de produtos cosméticos, peças de reposição e bens de consumo de baixa rotatividade, contexto direto desta dissertação \cite{Croston1972ForecastingIntermittentDemand,Boylan2008Intermittent}. Do ponto de vista gerencial, a intermitência torna políticas clássicas (como EOQ e $(s,S)$ parametrizadas pela média) subótimas, pois os pressupostos de demanda contínua e estacionária são violados sistematicamente.

Formalmente, seja $D_t \geq 0$ a demanda observada na loja no ciclo $t \in \{1,\ldots,T\}$. Define-se:

\begin{align}
  n_{+} &= |\{t : D_t > 0\}|, \label{eq:n_pos} \\
  \mu &= \frac{1}{T}\sum_{t=1}^{T} D_t, \quad \sigma^2 = \frac{1}{T-1}\sum_{t=1}^{T}(D_t - \mu)^2. \label{eq:mu_sigma}
\end{align}

O coeficiente de variação $\text{CV} = \sigma/\mu$ quantifica a dispersão relativa; valores $\text{CV} > 1$ indicam demanda altamente errática, característica do portfólio estudado (CV entre 1.5 e 5.4 no conjunto de dados empírico).

\subsection{Taxonomia de Syntetos-Boylan (ADI $\times$ CV$^2$)}
\label{sec:sb_taxa}

A classificação proposta por \citet{Syntetos2005Categorization} organiza séries de demanda em quatro quadrantes definidos por dois índices adimensionais:

\begin{align}
  \text{ADI} &= \frac{T}{n_{+}} \quad (\text{intervalo médio entre demandas, \emph{Average Demand Interval}}), \label{eq:adi} \\
  \text{CV}^2 &= \left(\frac{\sigma}{\mu}\right)^2 \quad (\text{quadrado do coeficiente de variação dos valores positivos}). \label{eq:cv2}
\end{align}

Os limiares empíricos $\text{ADI}^* = 1.32$ e $\text{CV}^{2*} = 0.49$, derivados analiticamente em \citet{Syntetos2005Categorization} a partir das condições de erro de Croston, definem quatro regimes:

\begin{table}[htbp]
\centering
\small
\renewcommand{\arraystretch}{1.2}
\caption{Quadrantes de Syntetos-Boylan: limiares e interpretações.}
\label{tab:sb_quad}
\begin{tabular}{llll}
\toprule
\textbf{Regime} & \textbf{ADI} & \textbf{CV$^2$} & \textbf{Interpretação operacional} \\
\midrule
\emph{Smooth} (Regular) & $< 1.32$ & $< 0.49$ & Alta frequência, demanda estável. \\
\emph{Erratic} (Errático) & $< 1.32$ & $\geq 0.49$ & Alta frequência, magnitude irregular. \\
\emph{Intermittent} (Intermitente) & $\geq 1.32$ & $< 0.49$ & Muitos zeros, magnitude relativamente estável. \\
\emph{Lumpy} (Grumoso) & $\geq 1.32$ & $\geq 0.49$ & Muitos zeros e magnitude altamente irregular. \\
\bottomrule
\end{tabular}
\end{table}

Essa taxonomia é amplamente adotada em literatura de supply chain \cite{Boylan2008Intermittent,Syntetos2012DemandClassificationsApplications} e constitui, nesta dissertação, a base para estratificação das análises: os KPIs das políticas são reportados por quadrante, permitindo identificar qual abordagem de controle é mais adequada a cada regime de demanda.

\subsection{Implicações para Política de Reposição}

A Figura~\ref{fig:sb_quadrant} ilustra a posição das séries no plano ADI $\times$ CV$^2$. Para o regime \emph{Smooth}, políticas baseadas em média (EOQ, Jornaleiro) são adequadas. Para regimes \emph{Intermittent} e \emph{Lumpy}, o método de Croston e suas variantes (SBA, TSB) são recomendados para previsão, enquanto políticas de reposição precisam ser calibradas por unidade para evitar ruptura sistêmica. Para o regime \emph{Erratic}, o maior risco é o excesso de estoque oriundo de pedidos baseados na média quando a demanda cai a zero.

\begin{figure}[htbp]
\centering
\begin{tikzpicture}[scale=1.1]
  \draw[->] (0,0) -- (5.2,0) node[right] {\small ADI};
  \draw[->] (0,0) -- (0,4.5) node[above] {\small CV$^2$};
  \draw[dashed,gray] (2.64,0) -- (2.64,4.2) node[above,font=\tiny] {$1.32$};
  \draw[dashed,gray] (0,1.96) -- (5,1.96) node[right,font=\tiny] {$0.49$};
  \node[align=center,font=\small] at (1.32,0.98) {\emph{Smooth}};
  \node[align=center,font=\small] at (1.32,3.10) {\emph{Erratic}};
  \node[align=center,font=\small] at (3.82,0.98) {\emph{Intermittent}};
  \node[align=center,font=\small] at (3.82,3.10) {\emph{Lumpy}};
\end{tikzpicture}
\caption{Plano ADI $\times$ CV$^2$ com os quatro quadrantes de Syntetos-Boylan. Limiares: ADI$^* = 1.32$, CV$^{2*} = 0.49$ \cite{Syntetos2005Categorization}.}
\label{fig:sb_quadrant}
\end{figure}

\section{Políticas de Controle de Inventário}

\subsection{Formulação do Problema de Controle}

O problema de controle de inventário em horizonte finito $T$ é formalizado como: dado o vetor de estado $\mathbf{s}_t = (I_t, Q_{t-L}, \hat{D}_{t+1})$ (nível de estoque corrente, pedidos em trânsito e previsão de demanda futura), determine uma política $\pi: \mathcal{S} \to \mathcal{A}$ que minimize o custo total esperado:

\begin{equation}
  \min_\pi \; \mathbb{E}\left[\sum_{t=1}^{T} \bigl(h\,I_t + c_s\,S_t + c_o^{\text{fix}}\,\mathbf{1}(Q_t>0) + c_o^{\text{var}}\,Q_t\bigr)\right],
  \label{eq:obj_inv_control}
\end{equation}

sujeito à dinâmica de estoque com prazo de entrega $L$:

\begin{equation}
  I_{t+1} = \max(0,\, I_t - D_t) + Q_{t-L}, \quad S_t = \max(0,\, D_t - I_t).
  \label{eq:balance}
\end{equation}

Os parâmetros $h$, $c_s$, $c_o^{\text{fix}}$, $c_o^{\text{var}}$ representam os custos unitários de manutenção, ruptura, fixo de ordenação e variável de ordenação, respectivamente. A razão $c_s/h = 5$ adotada nos experimentos reflete a perda de venda não realizada típica do varejo regional.

\subsection{Políticas Clássicas}

\paragraph{EOQ (\emph{Economic Order Quantity}).}
Minimiza o custo total anual assumindo demanda constante $\mu$. A quantidade ótima é $Q^* = \sqrt{2\mu c_o^{\text{fix}}/h}$, com ponto de reposição $s = \mu L$. Em regime intermitente, o pressuposto de demanda constante é violado, elevando o custo total em relação a políticas adaptativas.

\paragraph{Política $(s, S)$.}
Emite pedido de magnitude $S - I_t$ quando $I_t \leq s$, com $s$ e $S$ calibrados por EOQ. É a política de referência mais adotada no varejo regional brasileiro por sua simplicidade operacional, embora gere efeito \emph{bullwhip} elevado \cite{Lee1997BullwhipEffect}.

\paragraph{Jornaleiro (\emph{Newsvendor}).}
Determina o pedido ótimo de período único como o quantil $z_\alpha$ da distribuição de demanda: $Q^* = \hat{\mu} + z_\alpha \hat{\sigma}$, com $z_\alpha = 1.28$ para nível de serviço alvo de 90\%. Assume independência entre períodos, o que o torna econômico em demanda de baixo volume onde os picos não são correlacionados.

\subsection{Meta-Heurísticas de Otimização}

As meta-heurísticas buscam parâmetros $(s^*, Q^*)$ que minimizam o custo total via simulação do horizonte completo, sem pressupor qualquer distribuição de demanda. Essa abordagem é adequada ao regime intermitente porque integra diretamente os picos irregulares observados na série histórica.

\paragraph{Algoritmo Genético (GA).}
Mantém uma população de soluções e evolui por seleção, cruzamento e mutação \cite{Holland1992ANAS,Talbi2009Metaheuristics}. A função de aptidão é o custo total simulado. Em inventário, o GA tem sido amplamente utilizado para otimização de políticas de reposição \cite{Yuna2023IntermittentMetaheuristics}, pois tolera bem a descontinuidade introduzida pelos ciclos de demanda nula.

\paragraph{Recozimento Simulado (SA).}
Busca estocástica com critério de aceitação de Metropolis: aceita vizinho $s'$ com probabilidade $\exp(-\Delta C / T_k)$, onde $T_k$ decresce segundo esquema geométrico \cite{Kirkpatrick1983SA}. Permite escapar de ótimos locais mas é sensível ao esquema de resfriamento.

\paragraph{Otimização por Enxame de Partículas (PSO).}
Partículas navegam no espaço de soluções atraídas pelo melhor histórico pessoal e pelo melhor global, balanceando exploração e explotação via inércia $w$ e acelerações $c_1$, $c_2$ \cite{Kennedy1995PSO}.

\paragraph{Evolução Diferencial (DE).}
Gera candidatos por mutação diferencial (\emph{rand/1/bin}) e cruzamento com taxa $C_R$, sendo reconhecida pela robustez em espaços contínuos de baixa dimensão \cite{Storn1997DE}.

\subsection{Agentes de Aprendizado por Reforço}

O aprendizado por reforço (RL) enquadra o controle de inventário como um Processo de Decisão de Markov (MDP) $\langle \mathcal{S}, \mathcal{A}, P, R, \gamma \rangle$, onde o agente observa o estado $\mathbf{s}_t$, escolhe a ação $a_t$ (quantidade a pedir), recebe recompensa $r_t = -\text{custo}_t$ e busca maximizar o retorno descontado $\sum_{t}\gamma^t r_t$ \cite{SuttonBarto2018RL}. Essa formulação generaliza o controle ótimo sem impor pressuposto distribucional.

\paragraph{DQN.}
Aproxima a função-valor $Q_\theta(s, a)$ por redes neurais profundas, com \emph{experience replay} e \emph{target network} para estabilidade \cite{Mnih2015DQN}. Em inventário, exige espaço de ação discretizado e da ordem de $10^6$ interações para convergência, o que o torna inadequado quando a série histórica disponível é curta.

\paragraph{PPO.}
Algoritmo de \emph{policy gradient} com clipping $\epsilon$-clip que restringe a atualização da política, evitando colapsos de gradiente \cite{Schulman2017PPO}. Suporta ações contínuas (quantidades), tornando-o mais natural para controle de estoque. Sem inicialização adequada, tende a convergir para política degenerada de alta frequência de pedidos.

\paragraph{SARSA.}
Algoritmo tabular on-policy com atualização $Q(s,a) \leftarrow Q(s,a) + \alpha[r + \gamma Q(s',a') - Q(s,a)]$. Limitado em expressividade para regimes intermitentes por requerer discretização grosseira do espaço de estados.

\subsection{Arquitetura Híbrida GA-RL}

A arquitetura híbrida GA-RL, contribuição central desta dissertação, resolve a tensão entre convergência (RL) e horizontes curtos (série real com $T \approx 38$ ciclos). O pipeline opera em três fases:

\begin{enumerate}[(i)]
  \item \textbf{Busca global (GA)}: 50 gerações com 100 indivíduos geram 5.000 avaliações simuladas, identificando a região do espaço $(s, Q)$ com menor custo;
  \item \textbf{Pré-população do buffer}: as melhores soluções do GA populam o \emph{replay buffer} com 1.000 transições representativas de alta qualidade, eliminando a necessidade de exploração aleatória inicial;
  \item \textbf{Refinamento adaptativo (RL)}: o agente treina por 200 episódios a partir desse ponto qualificado, aprendendo a adaptar a política às flutuações específicas de cada loja.
\end{enumerate}

Essa arquitetura contorna o requisito de da ordem de $10^6$ passos do DQN isolado \cite{Oroojlooyjadid2022BeerGameDQN}, viabilizando políticas adaptativas mesmo com apenas 570 transições reais disponíveis ($15 \times 38$). A variante GA-PPO aplica a mesma lógica ao agente PPO contínuo, evitando a degeneração para $\text{FP} = 0.98$ observada no PPO isolado.

\section{Métricas de Qualidade de Previsão para Demanda Intermitente}

\subsection{Limitações das Métricas Clássicas}

As métricas clássicas MAE, RMSE e MAPE apresentam deficiências estruturais em regime intermitente: o MAPE é indefinido quando $D_t = 0$ (frequente na série estudada), o RMSE penaliza excessivamente erros em picos isolados, e nenhuma delas é invariante à escala, dificultando comparações entre séries de volume heterogêneo. A competição M4 \cite{Makridakis2020M4Competition} e \citet{Hyndman2006MASE} motivaram métricas especificamente desenvolvidas para contornar essas limitações.

\subsection{Métricas Adotadas}

Seja $\{D_t\}_{t=1}^{T}$ a demanda real e $\{\hat{D}_t\}_{t=1}^{T}$ a previsão. Definem-se os erros de previsão $e_t = D_t - \hat{D}_t$. A referência ingênua (\emph{naïve}) é o modelo de passeio aleatório: $\hat{D}_t^{\text{naïve}} = D_{t-1}$, com erro médio absoluto $\text{MAE}_{\text{naïve}} = \frac{1}{T-1}\sum_{t=2}^{T}|D_t - D_{t-1}|$.

\paragraph{MASE (\emph{Mean Absolute Scaled Error}).}
Proposto por \citet{Hyndman2006MASE}, é invariante à escala e bem definido mesmo com zeros na série:
\begin{equation}
  \text{MASE} = \frac{\text{MAE}}{\text{MAE}_{\text{naïve}}} = \frac{\frac{1}{T}\sum_t |e_t|}{\frac{1}{T-1}\sum_{t=2}^{T}|D_t - D_{t-1}|}.
  \label{eq:mase}
\end{equation}
$\text{MASE} < 1$ indica superação do modelo ingênuo; $\text{MASE} = 1$ é o \emph{benchmark} neutro.

\paragraph{sMAPE (\emph{Symmetric MAPE}).}
Proposto por \citet{Makridakis1993SMAPE}, utiliza denominador simétrico para lidar com zeros e assimetria:
\begin{equation}
  \text{sMAPE} = \frac{2}{T}\sum_t \frac{|e_t|}{|D_t| + |\hat{D}_t|} \times 100\%.
  \label{eq:smape}
\end{equation}
Limitado ao intervalo $[0\%, 200\%]$, sem singularidades quando $D_t = 0$ e $\hat{D}_t > 0$.

\paragraph{RMSSE (\emph{Root Mean Squared Scaled Error}).}
Versão escalada do RMSE, utilizada como métrica principal da competição M4 \cite{Makridakis2020M4Competition}:
\begin{equation}
  \text{RMSSE} = \sqrt{\frac{\frac{1}{T}\sum_t e_t^2}{\frac{1}{T-1}\sum_{t=2}^{T}(D_t - D_{t-1})^2}}.
  \label{eq:rmsse}
\end{equation}
Penaliza erros grandes (via quadrado) mas permanece comparável entre séries de diferentes escalas.

\paragraph{MBE (\emph{Mean Bias Error}).}
Detecta viés sistemático na previsão:
\begin{equation}
  \text{MBE} = \frac{1}{T}\sum_t e_t = \overline{D} - \overline{\hat{D}}.
  \label{eq:mbe}
\end{equation}
$\text{MBE} > 0$ indica subestimação sistemática (risco de ruptura); $\text{MBE} < 0$, superestimação (excesso de estoque).

\paragraph{U de Theil.}
Razão entre RMSE da previsão e RMSE do modelo ingênuo:
\begin{equation}
  U_{\text{Theil}} = \frac{\text{RMSE}}{\text{RMSE}_{\text{naïve}}}.
  \label{eq:theil}
\end{equation}
$U < 1$ indica superação do \emph{naïve}; valores $U \gg 1$ sinalizam modelo inadequado.

\subsection{Escolha da Métrica Principal}

O MASE é adotado como métrica principal nesta dissertação por três razões: (i)~invariância à escala permite comparar modelos entre lojas com volumes heterogêneos (razão de até 124× no conjunto de dados); (ii)~definição clara quando $D_t = 0$; (iii)~interpretação intuitiva em relação ao \emph{benchmark} ingênuo. O sMAPE é métrica secundária por sua bounded range e ampla adoção em benchmarks competitivos.

\section{Validação Temporal por Janela Deslizante (\emph{Walk-Forward})}
\label{sec:wf_protocol}

\subsection{Motivação}

A validação cruzada clássica (\emph{k-fold}) viola a ordem temporal das séries, introduzindo vazamento de informação futura no treinamento. Para séries temporais, o protocolo correto é a validação \emph{walk-forward} (ou \emph{rolling-origin}), que simula o processo de previsão online: o modelo é treinado exclusivamente em dados passados e avaliado em dados futuros, de forma sequencial.

\subsection{Protocolo Formal}

Dado horizonte total $T$ e janela mínima de treinamento $T_{\min}$, as divisões walk-forward são:

\begin{equation}
  \mathcal{D}_k = \bigl(\underbrace{\{1,\ldots,T_{\min}+k-1\}}_{\text{treino}},\; \underbrace{\{T_{\min}+k\}}_{\text{teste}}\bigr), \quad k = 1,\ldots,T - T_{\min}.
  \label{eq:wf_splits}
\end{equation}

O desempenho final é a média das métricas sobre todos os $T - T_{\min}$ pontos de teste. Essa abordagem garante que cada observação de teste foi prevista sem acesso a dados futuros, produzindo estimativas não viesadas da capacidade de generalização temporal do modelo.

\subsection{Implementação no Pipeline}

No pipeline Kedro desta dissertação, o nó \texttt{build\_walkforward\_splits} constrói as divisões \eqref{eq:wf_splits} com $T_{\min} = 12$ ciclos (mínimo para estimação de sazonalidade anual), produzindo até 26 pontos de teste por série. Os modelos --- Croston, ARIMA, XGBoost, LSTM --- são retro-alimentados a cada ciclo, e as métricas MASE, sMAPE, RMSSE são acumuladas por (loja, produto, modelo). Esse protocolo elimina o viés de avaliação \emph{in-sample} identificado como limitação central dos resultados da Fase~1.

\section{Testes Estatísticos Não Paramétricos para Comparação de Políticas}
\label{sec:stat_tests}

\subsection{Necessidade de Inferência Estatística}

A comparação entre 12 políticas sobre múltiplas lojas e múltiplos produtos na Fase~2 requer arcabouço inferencial formal. Como as distribuições de KPIs em regime intermitente são fortemente assimétricas, os testes paramétricos baseados em normalidade (ANOVA, $t$ de Student) são inadequados. Os testes não paramétricos (Wilcoxon, Friedman e Nemenyi) são adotados como padrão metodológico nesta dissertação.

\subsection{Teste de Wilcoxon com Sinal e Posto}

O teste de Wilcoxon \emph{signed-rank} \cite{Wilcoxon1945SignedRank} avalia se a mediana das diferenças pareadas entre duas políticas $\pi_A$ e $\pi_B$ é diferente de zero. Para o KPI $X$ (e.g., CTI) nas lojas $i = 1,\ldots,N$:

\begin{equation}
  W = \sum_{i : D_i \neq 0} \text{sgn}(D_i)\,R_i,
  \label{eq:wilcoxon}
\end{equation}

onde $D_i = X_i^A - X_i^B$ e $R_i$ é o posto de $|D_i|$. Hipótese nula: $\text{mediana}(D_i) = 0$. Ao contrário do $t$-test, não assume normalidade nem homocedasticidade.

\subsection{Teste de Friedman e \emph{Post-hoc} de Nemenyi}

Quando se compara $K > 2$ políticas simultaneamente, o teste de Friedman \cite{Friedman1940Ranks} testa se os postos médios de todas as políticas são iguais:

\begin{equation}
  \chi_F^2 = \frac{12N}{K(K+1)}\sum_{j=1}^{K}\left(\bar{r}_j - \frac{K+1}{2}\right)^2,
  \label{eq:friedman}
\end{equation}

onde $\bar{r}_j$ é o posto médio da política $j$ sobre as $N$ instâncias (lojas $\times$ produtos). A rejeição de $H_0$ indica que pelo menos uma política tem desempenho diferente. O \emph{post-hoc} de Nemenyi \cite{Nemenyi1963PostHoc} computa a diferença crítica (DC):

\begin{equation}
  \text{DC} = q_\alpha \sqrt{\frac{K(K+1)}{6N}},
  \label{eq:cd}
\end{equation}

onde $q_\alpha$ é o quantil da distribuição de Studentized range. Políticas cujos postos médios diferem em mais de DC são estatisticamente distinguíveis a nível $\alpha$.

\subsection{Tamanho de Efeito de Cohen}

O teste de Wilcoxon indica se a diferença é estatisticamente significante, mas não quantifica sua magnitude prática. O $d$ de Cohen \cite{Cohen1988StatisticalPower} fornece essa quantificação:

\begin{equation}
  d = \frac{\bar{X}_A - \bar{X}_B}{s_{\text{pooled}}}, \quad s_{\text{pooled}} = \sqrt{\frac{(N_A-1)s_A^2 + (N_B-1)s_B^2}{N_A+N_B-2}}.
  \label{eq:cohens_d}
\end{equation}

Convencionalmente: $|d| < 0.2$ é efeito negligenciável, $0.2 \leq |d| < 0.5$ pequeno, $0.5 \leq |d| < 0.8$ médio, $|d| \geq 0.8$ grande. Essa gradação é reportada para cada par (política proposta, baseline) em cada KPI e cada quadrante Syntetos-Boylan.

\section{Ambiente de Simulação e Formulação MDP}

O ambiente \emph{InventoryEnv}, implementado com API compatível com Gymnasium, simula o sistema de inventário $(s, Q)$ descrito pela Equação~\eqref{eq:balance}. O vetor de observação $\mathbf{o}_t \in \mathbb{R}^6$ normalizado contém: nível de estoque, demanda esperada, pedidos em trânsito, média móvel e desvio padrão dos últimos seis ciclos, e progresso do horizonte. A recompensa imediata é $r_t = -(h I_t + c_s S_t + c_o^{\text{fix}}\mathbf{1}(Q_t>0) + c_o^{\text{var}}Q_t)$, diretamente alinhada ao objetivo~\eqref{eq:obj_inv_control}.

A compatibilidade com Gymnasium permite que qualquer agente da biblioteca Stable-Baselines3 (DQN, PPO, A2C) treine no mesmo ambiente, garantindo comparação controlada entre arquiteturas de RL. A semente aleatória fixa (\texttt{seed=42}) e o uso de \emph{common random numbers} entre políticas eliminam o ruído amostral nas comparações.

\section{Arcabouço Tecnológico: Kedro e Reprodutibilidade}

\subsection{Kedro como Orquestrador de Pipeline}

Kedro \cite{KedroConcepts} estrutura o pipeline desta dissertação em cinco estágios conectados por um \emph{Data Catalog} declarativo (\texttt{catalog.yml}): (i)~\texttt{data\_ingestion}, (ii)~\texttt{demand\_forecasting}, (iii)~\texttt{inventory\_simulation}, (iv)~\texttt{statistical\_validation} e (v)~\texttt{reporting}. Cada estágio é composto por nós --- funções Python com entradas e saídas nomeadas --- cuja ordem de execução é inferida automaticamente pelo grafo de dependências.

Essa arquitetura garante três propriedades críticas para uma dissertação de mestrado:

\begin{itemize}
  \item \textbf{Rastreabilidade}: cada artefato (parquet, pickle, figura, tabela \LaTeX) tem origem determinística nos parâmetros YAML de entrada;
  \item \textbf{Reprodutibilidade}: \texttt{kedro run} reconstrói todos os artefatos do zero, independentemente da ordem de execução dos nós;
  \item \textbf{Configurabilidade}: alterar o escopo (e.g., de PB para todos os estados) requer apenas mudança no arquivo \texttt{conf/base/parameters/data\_ingestion.yml}, sem alteração do código.
\end{itemize}

\subsection{Estrutura do Catálogo de Dados}

Os dados são organizados em camadas numeradas segundo a convenção Kedro:

\begin{table}[htbp]
\centering
\small
\renewcommand{\arraystretch}{1.15}
\caption{Camadas do catálogo de dados do pipeline.}
\label{tab:catalog}
\begin{tabular}{lll}
\toprule
\textbf{Camada} & \textbf{Artefatos} & \textbf{Formato} \\
\midrule
\texttt{01\_raw} & Vendas particionadas por UF & Parquet (\emph{Hive}) \\
\texttt{02\_intermediate} & \texttt{sales\_raw}, \texttt{sales\_filtered}, \texttt{sales\_cleaned} & Parquet \\
\texttt{03\_primary} & \texttt{scenarios}, \texttt{scenarios\_meta} & Parquet \\
\texttt{04\_feature} & \texttt{walkforward\_splits} & Pickle \\
\texttt{06\_models} & \texttt{trained\_forecasters}, \texttt{scaled\_params} & Pickle \\
\texttt{07\_model\_output} & \texttt{forecast\_predictions}, \texttt{forecast\_metrics}, \texttt{kpis} & Parquet \\
\texttt{08\_reporting} & Figuras, tabelas \LaTeX, relatório & PDF/Tex/Pkl \\
\bottomrule
\end{tabular}
\end{table}

\subsection{Reprodutibilidade Experimental}

A reprodutibilidade do pipeline é garantida por: (i)~versionamento Git de código e configurações YAML; (ii)~sementes globais fixas por experimento (\texttt{random\_seed: 42}); (iii)~\emph{common random numbers} para simulações comparativas entre políticas; (iv)~hashes de datasets registrados no catálogo; e (v)~containerização Docker com dependências fixadas no \texttt{requirements.txt}. Esses procedimentos seguem as diretrizes de reprodutibilidade experimental em pesquisa computacional \cite{Pineau2021ReproducibilityML}.
